\documentclass[../main.tex]{subfiles}
\begin{document}
%==========================================TIÊU ĐỀ TẬP========================================
\begin{table}[h]
    \begin{adjustwidth}{-1cm}{}
        \begin{tabular}{>{\centering\arraybackslash}p{6cm}|>{\raggedright\arraybackslash}p{12.5cm}}
            \multirow{5}{6cm}
            {\\[-40pt]
                \flaregothic\fontsize{20pt}{20pt}\selectfont \centering
                \color{eptype}HISTOIRE\color{black}\\
                \freshbot\fontsize{72pt}{72pt}\selectfont
                \color{epnum}25\\[6pt]
                \cafeteria\fontsize{20pt}{20pt}\selectfont\color{epnum}(D-25)\color{black}\\[10pt]
                \sen\fontsize{18pt}{18pt}\selectfont
                \color{epattr}Mini-histoire
                \color{black}
            }
             &
            {\fotskip\fontsize{18pt}{18pt}\selectfont \ruby{五}{ご}\ruby{色}{しき}のholoFive}
            \\[6pt]
             & {\cafeteria\fontsize{18pt}{18pt}\selectfont Five Colors of holoFive}       \\[4pt]
             & {\vnmsans\fontsize{18pt}{18pt}\selectfont Năm sắc màu holoFive}            \\[8pt]
             & {\caviar{\fontsize{14pt}{14pt}\selectfont Nhân vật: Theo từng câu chuyện}} \\[8pt]
        \end{tabular}
    \end{adjustwidth}
\end{table}
%=========================================TRUYỆN CHÍNH========================================
\color{black}

\txtbx[2]{
    {\color{vol} \uline{Tóm tắt:}} Vài tuần sau ngày hội ngộ, Delutaya lần lượt có những buổi gặp riêng với từng thành viên NePoLaBo. Một tách trà với Lamy mở ra câu chuyện về nỗi sợ cô đơn; buổi nấu ăn với Nene biến căn bếp thành bãi chiến trường; đêm chơi game kinh dị cùng Polka khiến Kaigou phải can thiệp; lời khuyên bình tĩnh của Botan giúp Delutaya nhìn nhận sự mệt mỏi khác đi. Cuối cùng, cả năm người cùng hát karaoke trong phòng thu, khép lại chuỗi ngày bằng một bài hát mới của Delutaya.
}

\pov{Delutaya}

\stpt{-Histoire 1-}
\stcap{Ly trà chiều ấm áp}
\begin{center}
    {\caviar\fontsize{12pt}{12pt}\selectfont Nhân vật: \emp{delu} \emp{lamy}}
\end{center}

Vài tuần trôi qua kể từ buổi hội ngộ của \textit{holoFive} tại căn nhà gia đình Hikawa, nhịp sống thường nhật của tôi vẫn phẳng lặng trôi như một thói quen. Sáng nào tôi cũng thức dậy đúng giờ, tự tay đun nước pha ấm trà nóng, lướt qua hàng dài thông báo công việc rồi lại nhốt mình cả ngày trong phòng thu cùng những giai điệu còn dang dở.

Chiếc điện thoại trước đây vốn lặng lẽ nằm một góc, nay thỉnh thoảng lại sáng lên bởi những dòng tin nhắn rủ rê đi ăn, dạo phố, hay chỉ đơn giản là tạt qua chơi. Và người mở lời đầu tiên là Lamy.

\txtbx{``Delu-chan, chiều nay bà có nhà không? Tôi sang chơi nhé\~{}''}

Tôi đọc đi đọc lại câu thoại ấy tới ba lần, cứ ngỡ mình nhìn nhầm tên người gửi. Sau khi gõ câu trả lời ngắn gọn rằng tôi ở nhà cả ngày rồi bấm gửi, tôi vẫn giương mắt nhìn trối trối vào màn hình. Dòng chữ ấy vẫn nằm nguyên ở đó, vừa chân thực vừa dịu dàng như một lời hẹn hết sức bình thường.

Tôi giật mình ngước nhìn căn phòng thu chất đống những tờ giấy ghi chú cùng vô số bản nhạc in chưa kịp dọn. Chợt nghĩ mình nên tút tát lại không gian một chút, tôi vội vã gom gọn chồng tài liệu trên bàn. Vừa đặt trang giấy cuối cùng xuống thì tiếng chuông cửa vang lên.

Lamy đứng trước hiên nhà, mái tóc được buộc gọn gàng đằng sau, tay ôm một chiếc túi vải nhỏ nhắn: ``Hóa ra bà có nhà thật này.''

``Ừm, vào nhà đi, Lamy-chan.''

Bà ấy tháo đôi giày thể thao đặt gọn gàng ngoài cửa, bước lên phòng khách tầng hai rồi chìa chiếc túi vải về phía tôi. ``Tôi có mua ít bánh quy ở tiệm gần nhà. Không biết bà thích vị gì nên tôi lấy cả socola lẫn matcha luôn.''

Tôi đón lấy túi bánh, lồng ngực khẽ dâng lên một luồng ấm áp trước sự chu đáo dịu dàng của bà ấy. ``Cảm ơn bà nhé. Tôi cũng đang định đi pha trà đây.''

``Thế để tôi phụ bà một tay.''

Cả hai cùng bước vào bếp. Tôi đi đun ấm nước sôi, còn bà ấy tự nhiên như thể ở nhà mình, mở tủ chén lấy ra hai chiếc cốc sứ trắng mà tôi mới sắm tháng trước. Bà hỏi tôi muốn uống trà đen hay trà thảo mộc, tôi bảo thích trà hoa cúc, loại trà quen thuộc mà cả hai từng uống rất nhiều hồi còn ở chung nhà.

Lamy nhẹ nhàng mở gói bánh, xếp từng chiếc gọn gàng thành hai hàng trên đĩa nhỏ. Thấy bà ấy chăm chú rót làn nước nóng ngạt ngào hương hoa vào cốc, tôi không nén nổi nụ cười, bèn nhíu mày hỏi trêu: ``Này Lamy-chan, hôm nay trời đổ mưa ngược hay sao mà ``tửu thần' nhà ta lại chuyển sang thanh tịnh uống trà thảo mộc thế này? Tôi còn tưởng vừa mở cửa ra là bà đã tay xách chai sake chứ?''

Lamy khựng lại một nhịp, đôi má hơi ửng hồng rồi bĩu môi lườm tôi một cái đầy giận dỗi: ``Này nhé, Delu! Tôi đến thăm bà đàng hoàng đấy nhé. Chẳng lẽ sang nhà bạn thân chơi ban ngày ban mặt mà tôi lại vác chai rượu ra tu ừng ực sao? Với lại\dots\ uống trà mới hợp vị bánh quy matcha này chứ! Lâu lâu cũng phải cho gan nghỉ ngơi chút đỉnh chứ bộ!''

Nhìn điệu bộ phân bua cuống quít ấy, tôi bật cười thành tiếng. Sự e ngại hay khoảng cách dường như đã bốc hơi sạch sẽ theo làn khói trà hoa cúc nghi ngút.

Lamy thao tác rất chậm rãi, không hối hả gợi lại buổi hội ngộ vừa rồi, cũng chẳng mổ xẻ những tháng ngày xa cách đã qua. Mấy câu trả lời mà tôi dày công chuẩn bị sẵn trong đầu hóa ra lại trở thành thừa thãi.

\ct

Chúng tôi kéo ghế ngồi bên khung cửa sổ lớn hướng ra khoảng trời xanh ngắt. Lamy khẽ thổi cho làn hơi nước tan bớt rồi nhấp một ngụm trà nhỏ.

``Dạo này công việc của bà có bận lắm không?'' bà lên tiếng trước.

``Cũng bình thường thôi. Tôi đang tập trung hoàn thiện mấy bản phối mới nên đôi khi hay thức khuya một chút.''

``Thế có ngủ nghê tử tế không đấy?''

Tôi khựng lại một chút rồi thành thật: ``Hôm nào ngoan thì ngủ đủ, còn hôm nào cuốn theo giai điệu quá thì quên mất giờ giấc, ngẩng đầu lên đã thấy bình minh rồi.''

Lamy mỉm cười nhẹ nhàng, nhưng ánh mắt vàng lấp lánh sự lo lắng chân thành. Bà đặt cốc trà xuống, rồi nói: ``Ngày xưa tôi cũng từng như thế. Có khoảng thời gian, tôi luôn sống trong cảm giác sợ hãi rằng nếu mình chỉ cần khựng lại dù chỉ một bước thôi, mọi người sẽ dần quên mất sự tồn tại của mình.''

Tôi đặt cốc trà xuống mặt bàn gỗ, lồng ngực hẫng đi một nhịp khi nghe những lời đó. Lamy nhìn ra khoảng không ngoài cửa sổ, tông giọng trầm hẳn xuống: ``Từ ngày Aloe-chan rời đi\dots\ tôi đã sợ lắm. Tôi sợ năm đứa rồi cũng sẽ xa cách. Sợ một ngày nào đó mỗi người một ngả, cuốn theo guồng quay riêng đến mức quên cả số điện thoại của nhau, chẳng còn lý do để cất lời hỏi thăm một câu đơn giản\dots''

Tôi lặng người, không biết nên đáp lại làm sao. Lamy không nhìn tôi, bà ấy không cần những lời an ủi sáo rỗng hay thương hại. Bà chỉ đang trải lòng về những vết thương giấu kín suốt nhiều năm qua.

``Tôi cũng từng sợ điều đó lắm\dots'' tôi thì thào, giọng nghẹn lại.

``Thế còn bây giờ thì sao?''

Tôi liếc nhìn chiếc điện thoại nằm im bên cốc trà, rồi ngước mắt nhìn người bạn đối diện: ``Bây giờ\dots\ bà đang ngồi ở đây với tôi rồi.''

Lamy ngẩn ra một giây rồi nở nụ cười rạng rỡ, dịu dàng như vệt nắng chiều rọi qua khung cửa. ``Ừ. Tôi đang ở đây rồi.''

\ct

Sau giây phút lắng đọng ấy, không khí lại trở về với những câu chuyện phiếm không đầu không cuối như hồi chúng tôi chưa rời xa. Lamy kể về buổi livestream gần đây, khi bà vô tình gạt tay làm vỡ tan cốc nước ngay trên sóng. Tôi cũng phì cười đáp lại bằng sự cố ngớ ngẩn của chính mình: dạo này hay quên lưu file dự án, đến lúc máy tính dở chứng ngắt đột ngột mới hoảng hốt đổ mồ hôi hột.

``Bà đã cài ứng dụng ghi chú nhắc nhở vào điện thoại chưa đấy?'' Lamy hỏi.

``Tôi cài hẳn ba cái rồi, nhưng vẫn quên như thường. Vì căn bản tôi có bao giờ chịu mở mấy cái app đó ra đâu,'' tôi lắc đầu bất lực.

Lamy bật cười thành tiếng: ``Vậy thì bà phải viết chữ thật to lên một tờ giấy rực rỡ, rồi dán chình ình ngay giữa màn hình máy tính ấy!''

``Thế thì tôi cũng sẽ quên đọc nó thôi,'' tôi tặc lưỡi.

``Vậy thì dán ngay lên cánh cửa phòng! Ra vào kiểu gì chẳng phải nhìn thấy?''

``Thế nhỡ tôi quên luôn cả việc mở cửa phòng thì sao?''

Cô bạn Yêu tinh tuyết liền nhăn mặt: ``Này nhé Delu-chan, bà đang cố tình bắt bẻ tôi đấy à?''

Tôi bật cười nghiêng ngả trước gương mặt giả vờ hờn dỗi nhưng đôi mắt vẫn tràn ngập ý cười của Lamy. Bà ấy thong thả lấy thêm một chiếc bánh quy socola, nhai nhè nhẹ. Buổi chiều trôi qua dịu dàng và thanh bình đến lạ. Những khoảng lặng giữa hai chúng tôi chẳng hề gượng gạo, mà êm đềm như một lẽ tự nhiên.

Trước khi ra về, Lamy nghiêng đầu đề nghị: ``Nếu bà không phiền, thì tuần nào chúng mình cũng dành ra một buổi chiều để cùng uống trà như thế này nhé?''

Không một giây chần chừ, tôi gật đầu ngay: ``Ừm, được chứ.''

Bà ấy không đưa ra lời hứa hẹn đao to búa lớn rằng năm đứa chúng tôi sẽ mãi mãi bên nhau không chia lìa. Bà chỉ nhẹ nhàng hẹn gặp lại vào tuần sau, và những tuần kế tiếp.

Ngay khoảnh khắc ấy, tôi chợt nhận ra rằng, đôi khi chỉ cần một lời hẹn giản đơn như thế thôi, cũng đủ để xua tan đi những nỗi sợ vô hình tích tụ bấy lâu trong lòng.

\newpage

\stpt{-Histoire 2-}
\stcap{Đầu bếp hỗn loạn}
\begin{center}
    {\caviar\fontsize{12pt}{12pt}\selectfont Nhân vật: \emp{delu} \emp{nene} \emp{suzuno2}}
\end{center}

Buổi gặp thứ hai bắt đầu bằng một tin nhắn thoại dài gần một phút. Giọng của Nene-chi vang lên đầy phấn khích qua loa điện thoại: ``\textbf{Delu-chan! Hôm nay tôi dạy bà nấu một món thật ngon nha\~{}!}''

Tôi nghe lại tin nhắn, rồi nhìn chiếc tạp dề đang cầm trên tay. Thực sự mà nói, tôi không giỏi nấu ăn, ngoài món cơm cuộn sushi mà tôi và chị Mikeneko đã thi đấu từ vài tháng trước. Thêm nữa, ký ức mà anh Kuon đánh giá món ăn tôi làm từ hồi tôi mới dọn tới ở vẫn chưa hề nguôi ngoai.

Tôi cũng không biết Nene giỏi đến đâu; hồi tôi còn sống chung với mọi người, hiếm khi tôi nhìn thấy cô ấy vào bếp.

Dù trong lòng còn chút lo lắng về kỹ năng nấu ăn của bản thân, nhưng sự hào hứng trong giọng nói của Nene quá lây lan, nên cuối cùng tôi cũng gật đầu đồng ý với lời mời. Tôi nhắn lại với cô ấy rằng mọi thứ đều ổn, và chúng ta có thể gặp nhau vào cuối tuần tại căn bếp chung của nhà.

Sau khi cúp máy với Nene, tôi bước ra khỏi phòng thu, tình cờ gặp Suzuno đang lau dọn kệ sách ở hành lang. Em ấy ngước nhìn tôi và mỉm cười: ``Chuyện gì vui mà chị cười tươi thế ạ, Delutaya-san?''

Tôi vừa đi đến bên em vừa kể hết câu chuyện: ``À, Nene-chi mới rủ chị học nấu cơm rang vào cuối tuần này. Chị không chắc là gian bếp ở tầng hai của chị có đủ rộng không nữa\dots''

Em ấy lắng nghe xong, đôi mắt sáng lên rồi liền xua tay: ``Ồ, không cần đâu chị ạ. Căn bếp chung ở tầng sáu rộng rãi và có đủ tất cả dụng cụ rồi. Để em sắp xếp cho hai chị dùng vào cuối tuần nhé. Em còn có thể chuẩn bị sẵn một ít nguyên liệu cơ bản như gạo, trứng, rau củ để hai chị không phải mua đâu.''

``Tiện quá vậy! Em có chắc không làm phiền công việc của em chứ?'' tôi hỏi, ngại ngùng vì lại làm thêm việc cho em ấy.

Suzuno liền cười lớn, lắc đầu nguầy nguậy: ``Không có gì phiền đâu ạ! Em cũng muốn xem hai chị Nene-san và chị nấu ăn thế nào nữa. Hứa sẽ không làm phiền hai chị đâu, chỉ đứng quan sát thôi ạ.''

Thế là em ấy nhanh chóng sắp xếp lịch để hai người chúng tôi có thể dùng thoải mái căn bếp rộng rãi ở tầng ba. Em ấy còn cẩn thận chuẩn bị sẵn tất cả nguyên liệu cần thiết, phân loại chúng gọn gàng vào những chiếc hộp nhỏ trên bàn bếp, để khi chúng tôi đến chỉ cần lấy ra dùng là được.

``Em chuẩn bị sẵn một số nguyên liệu cơ bản rồi ạ. Nếu Nene-san và Delutaya-san cần thêm gì thì cứ gọi em nhé,'' Suzuno mỉm cười lễ phép nói.

\ct

Khoảng hai mươi phút sau, Nene tới nhà, mang theo một cuốn sổ công thức dày cộp và chiếc băng đô màu cam nổi bật trên đầu. Cô ấy lật mở cuốn sổ, bên trong chật kín những hình vẽ phác họa ngốc nghếch cùng mấy dòng chữ nguệch ngoạc.

``Hôm nay mình sẽ cùng làm cơm rang nha! Món này làm cực kỳ đơn giản thôi!''

Tôi lướt mắt qua từng trang sổ, lòng thầm ngạc nhiên trước đống hình vẽ hoạt hình ngộ nghĩnh chất đầy giấy, đôi chân vẫn đứng yên cạnh bếp như thể đang chờ đợi điều gì đó sắp xảy ra.

``Bà chắc chứ? Nó thật sự dễ làm hả?'' tôi nhíu mày, giọng pha chút nghi ngờ khi còn nhớ lần nấu ăn thất bại của mình cách đây không lâu.

``Tôi làm món này cả trăm lần rồi, bà hoàn toàn có thể tin tưởng vào khả năng của tôi!'' Nene hơi nhún người về phía trước, một tay vỗ mạnh vào ngực, trông tự tin như thể đây là đầu bếp hàng đầu của một nhà hàng lớn.

Suzuno đứng góc bếp, hai tay đặt nhẹ trước bụng, mắt dõi theo từng cử động của chúng tôi. Em ấy khẽ nghiêng người, giọng nhẹ nhàng như sợ làm phiền không khí hào hứng: ``Em đã chuẩn bị sẵn hết tất cả dụng cụ ở đây cho hai chị rồi ạ.''

Cô bạn tóc vàng lập tức kéo chiếc ghế nhựa gần đó lại sát bếp, bước lên đứng ổn định, sau đó xắn gọn hai tay áo lên đến khuỷu tay, tư thế trông chẳng khác nào một chuyên gia ẩm thực chuẩn bị thể hiện kỹ năng.

``Bước đầu tiên, mình cần làm nóng chảo trước đã.'' Cô ấy chỉ vào chiếc chảo chống dính đặt trên bếp, giọng đầy chỉ huy. Tôi với tay bật công tắc bếp ga. Một ngọn lửa xanh lập lòe bùng lên dưới đáy chảo.

Nene cúi gập cả người xuống, mặt gần như chạm vào mặt bếp, mắt chăm chú nhìn ngọn lửa như đang kiểm tra một thiết bị quan trọng: ``Lửa hơi lớn quá rồi, nhỏ lại một chút đi.''

Tôi nghe lời, đưa tay vặn núm điều chỉnh lửa sang bên trái thêm một nấc. Đột nhiên, ngọn lửa tắt hẳn. Căn bếp chìm vào im lặng trong vài giây.

Chúng tôi cứng người nhìn nhau, mắt đều mở to bất ngờ. Nene chớp chớp mắt vài lần như thể vừa xử lý một tình huống bất ngờ, rồi bỗng gật gù cái đầu, gương mặt thản nhiên như thể mọi thứ đều nằm trong kế hoạch của mình: ``Ồ, cái bếp này hiện đại thật đấy, tự biết tắt để tiết kiệm gas luôn!''

Suzuno đứng cạnh chỉ có thể khẽ thở dài một tiếng, tiếng thở dài chứa đựng sự bất lực trước cô bạn dù lớn tuổi vẫn còn ngốc nghếch. Em ấy bước lại, nhẹ nhàng đưa tay bật lại bếp, ngọn lửa xanh lại lập lòe bùng lên: ``Để em bật lại bếp giúp hai chị ạ.''

Nene chỉnh chiếc băng đô màu cam trên trán, rồi chỉ vào củ hành tây còn nguyên trên thớt: ``Bước tiếp theo mình cắt hành thôi! Để tôi chỉ bà từng bước một.''

Tôi gật đầu vội, hai tay cầm chặt con dao làm bếp, điều chỉnh tư thế cho đúng theo mấy lần thấy người khác nấu ăn. Tôi thở một hơi nhẹ rồi đưa lưỡi dao xuống củ hành.

Miếng hành đầu tiên rời thân củ, nó hơi dày cộp gần bằng nửa ngón tay. Tôi cau mày chút xíu, lưỡng lự cắt miếng tiếp theo, nhưng miếng này lại mỏng gần như trong suốt. Cố gắng điều chỉnh lực tay, miếng thứ ba lại bị cắt dài ngoẵng, gấp đôi chiều dài của hai miếng trước đó. Tôi ngồi tròn mắt nhìn đống hành bất đối xứng trên thớt, trong bụng bắt đầu lo lắng mình lại phá hỏng nguyên liệu.

Nene cũng dán mắt vào đống hành ``độc đáo'', hai lông mày khẽ cau lại suy ngẫm mấy giây. Bỗng cô ấy phất tay hào phóng, mặt nở nụ cười tươi rói như vừa phát hiện ra một trường phái ẩm thực mới: ``Không sao đâu! Món cơm rang này là phong cách tự do, không cần các miếng phải đều tăm tắp như nhà hàng đâu. Thế này mới có hương vị nhà làm!''

Tôi thở phào nhẹ nhõm, vai thả lỏng ra sau khi nghe lời an ủi. Đúng là chỉ có Nene mới có thể biến sự vụng về của tôi thành một ``phong cách'' độc đáo.

Rồi chúng tôi chuyển sang bước tiếp theo: đập trứng vào chiếc bát sứ lớn ở giữa bàn. Tôi cầm một quả trứng, nghiêng nhẹ mép bát để đập vỡ vỏ, nhưng có lẽ tôi hơi mạnh tay quá. Một mảng vỏ trứng to đùng rơi tọt thẳng vào lòng đỏ đang nằm yên trong bát, tôi đứng hình nhìn mảnh vỏ trôi nổi trên mặt trứng.

Cô bạn lập tức nghiêm mặt, hai vai thẳng lên như đang chuẩn bị thực hiện một nhiệm vụ tối mật. Cô ấy nhón hai ngón tay trỏ và cái, chậm rãi, cẩn thận đến từng li từng tí gắp mảnh vỏ trứng ra khỏi bát, phong thái trang trọng không khác gì đang nâng niu một bảo vật quốc gia.

Đặt mảnh vỏ vào góc thớt, Nene vẫn giữ vẻ mặt nghiêm chỉnh, quay sang tôi nói với giọng đều đều đến buồn cười: ``Đây mới là phần khó nhất của mọi công việc nấu ăn đấy.''

Tôi nhíu mày thắc mắc, chỉ vào chiếc bát trứng: ``Là gắp vỏ trứng ra hả?''

``Không phải,'' Nene lắc đầu, vẫn không thay đổi sắc mặt. ``Là giữ được vẻ mặt ngầu để người khác không phát hiện ra mình đang hoảng sợ thực sự.''

Cả căn bếp bỗng dưng vang lên tiếng cười của tôi, tôi cười nghiêng ngả, phải vịn vào cạnh bàn để không ngồi sụt xuống đất. Suzuno đứng ở góc bếp cũng không nhịn được, bụm miệng cười rồi bước tới đưa cho chúng tôi một chiếc bát sạch khác, để chúng tôi có nơi đựng những mảnh vỏ trứng vụn vặt.

Sau khi giải quyết xong đống vỏ trứng, Nene cầm chiếc phới lồng bắt đầu đánh đều hỗn hợp trứng. Phần việc của tôi là đo nước tương đổ vào cơm sau này. Tôi cúi đầu đọc dòng chữ nguệch ngoạc trong cuốn sổ công thức của Nene: ``Công thức ghi cần một thìa cà phê nước tương.''

Vội vàng vươn tay ra với hộp đựng thìa, tôi lại cầm nhầm chiếc thìa múc canh lớn bềnh bồng. Nene đứng bên cạnh không những không nhắc ra, còn nhiệt tình gật đầu lia lịa, vỗ tay cổ vũ: ``Đúng rồi đó! Múc đầy thìa đó luôn cho đậm đà, món ăn mới ngon!''

Suzuno đứng gần đó vừa rửa tay nghe thấy thế, liền vội vàng bước tới, nhẹ nhàng giữ lấy cổ tay tôi trước khi nước tương kịp chảy một giọt nào ra khỏi chai. Em ấy cười khúc khích, giọng nói vẫn lễ phép: ``Delutaya-san, thìa đó là thìa canh lớn rồi ạ. Thìa cà phê nhỏ đúng công thức là chiếc này cơ ạ.'' Em ấy chỉ vào chiếc thìa nhỏ xíu nằm góc hộp.

Tôi bừng tỉnh, rụt tay lại cầm đúng chiếc thìa cà phê, mặt hơi ửng hồng vì ngượng: ``À\dots\ cảm ơn Suzuno-san nhé, nếu không có em thì có lẽ cả nồi cơm rang đã thành hồ nước tương mất.''

Nene đứng bên cạnh cũng gãi đầu gãi tai, cười ngượng ngùng không kém, hai má hơi đỏ lên: ``Haha\~{}\dots\ T-Tôi cũng đang định nhắc bà thế đấy! Có vẻ hôm nay đầu bếp Nene cũng cần một người giám sát riêng rồi, để không đổ cả lọ gia vị vào món ăn lại thành dở.''

\ct

Sau vài phút, nhiệt độ trong chảo đã đủ nóng. Nene nghiêng chai dầu, để một lượng nhỏ dầu ăn chảy loang đều khắp mặt chảo, sau đó bà ấy quay sang giật giật tay áo tôi, mắt sáng lên như vừa đưa ra mệnh lệnh chiến dịch quan trọng nhất: ``Đổ trứng đi nào!''

Tôi hít một hơi, nâng chiếc bát đựng trứng đã đánh đều lên, nghiêng nhẹ để hỗn hợp vàng óng chảy vào chảo. Ngay khi trứng chạm vào mặt kim loại nóng, một tiếng ``xèo'' lớn vang lên như tiếng pháo hoa nổ, kèm theo một làn khói trắng bay thẳng lên trần nhà. Tôi giật mình lùi lại nửa bước, tim đập nhanh như vừa bị một con chuột nhảy ra từ góc bếp.

Nene thì phản ứng ngược lại, nàng lập tức chộp lấy đôi đũa gỗ, lao tới trước chảo và đảo trứng một cách nhiệt tình, tóc mai hơi bay theo động tác dũng mãnh: ``Đừng sợ tiếng đó! Nó chỉ đang gào lên đòi chín thôi!''

Tôi gật đầu lia lịa, tay vẫn còn hơi run khi cầm lấy chiếc đĩa lớn đựng cơm nguội. Nene vừa đảo trứng vừa quay sang hét nhỏ với tôi: ``\textbf{Đổ cơm vào đi nào!}''

Tôi cố gắng bình tĩnh, nghiêng chiếc đĩa thật đều, nhưng chỉ có khoảng nửa số cơm rơi gọn gàng vào chảo. Nửa còn lại thì như một đoàn quân nhảy dù, bay tự do trên không trung và đáp thẳng xuống mặt bếp đá hoa rồi rơi bẹp xuống, tạo ra một đống cơm tròn tròn ngay trước chân tôi.

Cả hai chúng tôi lập tức đứng hình, mắt tròn mắt dẹt nhìn đống cơm trên sàn. Nene chớp chớp mắt vài lần, rồi đột nhiên bật cười khúc khắc, tiếng cười ngày càng lớn dần, cuối cùng nàng phải ôm bụng và vịn vào cạnh bàn để không ngồi sụt xuống: ``Haha! Cơm biết nhảy dù kìa! Tôi xin lỗi\dots\ tôi không nên cười\dots\ nhưng trông nó buồn cười quá!''

Suzuno từ góc bếp nhanh nhẹn bước tới, tay mang cây vét bột và một chiếc đĩa sạch, em ấy mỉm cười dịu dàng khi nhìn thấy cảnh tượng: ``Không sao đâu ạ, mình nhặt phần sạch trên mặt bếp trước nhé hai chị.''

Tôi cũng bật cười theo, nhìn Nene rồi lại nhìn đống cơm vô tội: ``Đúng vậy, nếu bà không cười thì tôi cũng chẳng biết phải phản ứng làm sao lúc này nữa.''

Chúng tôi gom được phần cơm còn sạch, bỏ vào chảo và tiếp tục đảo đều hỗn hợp. Nene đổ đống hành tây, trứng và mớ rau củ tôi vừa cắt xong (dù chúng vẫn còn hơi to nhỏ không đều) vào, còn tôi thì bắt đầu nêm gia vị từng chút một.

Suzuno đứng bên cạnh nhắc nhở rất chu đáo, giọng nói vẫn như một người mẹ chu toàn như mọi khi: ``Hai chị nhớ nếm thử trước khi cho thêm gia vị ạ.''

Lần đầu tiên chúng tôi múc một ít ra nếm, món ăn hơi nhạt so với khẩu vị. Nene lập tức tỏ ra rất tự tin, múc hẳn một thìa muối đầy và đổ thẳng vào chảo, tay còn vung vẩy tán dương: ``Phải thế này mới chuẩn vị!''

Nhưng lần nếm thứ hai\dots\ cả hai tôi đều nhăn mặt lại, vị mặn chát lan tỏa khắp miệng. Cô bạn trợn tròn mắt, lập tức chạy tới lấy nồi cơm và đổ thêm một lượng lớn cơm nguội vào chảo, như thể đang cố gắng ``chữa cháy'' cho món ăn: ``Không sao đâu! Thêm thật nhiều cơm vào là hòa tan hết mặn ngay!''

Nhưng chúng tôi vừa mới đảo cơm được vài vòng, mùi khói bốc lên mỗi lúc một đậm, làm cho chuông báo khói trên trần nhà kêu lên một tiếng *bíp* ngắn, đủ làm cả ba người trong bếp nhảy lên. Tôi giật nảy mình, suýt đánh rơi chiếc thìa đang cầm trên tay. Nene thì hốt hoảng tắt vội bếp ga, sau đó cuống quýt dùng hai tay quạt khói một cách vô vọng, như thể muốn thổi bay hết mọi dấu vết của vụ cháy nhỏ.

Suzuno vẫn giữ được sự bình tĩnh hiếm có, em ấy nhanh chóng mở tung tất cả cửa sổ trong bếp, rồi lấy chiếc khăn sạch quạt nhẹ về phía cảm biến khói, vừa làm vừa nói để trấn an chúng tôi: ``Không sao đâu ạ. Chỉ là chút hơi nóng chạm vào cảm biến thôi, hai chị yên tâm nhé.''

Tôi thở phào nhẹ nhõm, rồi nhìn vào chảo cơm. Một góc dưới đáy chảo đã chuyển sang màu đen sém, trông như vừa trải qua một trận chiến nhỏ. Nene cũng nhìn thấy điều đó, nhưng thay vì hoảng sợ, nàng lại nhìn tôi với đôi mắt sáng rực, lấy lại sự tự tin một cách kỳ lạ: ``Tôi quyết định rồi! Món này tên là 'Cơm rang giòn cháy cạnh đặc sản'!''

Tôi nhướng mày, không thể tin vào lời nàng nói: ``Tên đó có thật không đấy?''

``Bây giờ thì có rồi đó! Do chính siêu đầu bếp Nene-chi này vừa sáng tạo ra!'' cô bạn đập ngực tự hào.

Suzuno mỉm cười tiếp lời, khéo léo giúp chúng tôi xúc phần cơm không bị cháy sang chiếc đĩa lớn: ``Dạ, em chuyển phần cơm ngon sang đĩa cho hai chị nhé.''

Nene hân hoan rắc một nhúm hành lá tươi lên trên đĩa cơm, còn tôi thì đặt hai chiếc thìa cạnh đĩa. Cả ba chúng tôi cùng nhau nếm thử thành quả của mình.

Vị cơm hơi mặn, rau củ thì miếng to miếng nhỏ, vài miếng còn hơi cứng, cộng thêm chút mùi khói đặc trưng từ vụ cháy nhỏ. Nhưng kỳ lạ thay, nó vẫn ăn được, thậm chí còn có một hương vị ấm áp khó tả.

Cô bạn nhai chậm rãi, nét mặt biểu cảm đầy hài hước khi nhìn vào đĩa thức ăn: ``Ừm\dots\ phải thừa nhận là nó không phải tuyệt tác đỉnh cao nhất của tôi.''

Tôi gật đầu đồng tình, nuốt hết miếng cơm trong miệng: ``Nhưng cũng không đến mức tệ hại đâu.''

Suzuno cũng lễ phép nếm thử một miếng nhỏ từ góc đĩa, em ấy gật đầu rồi góp ý một cách nhẹ nhàng: ``Lần sau nếu cắt rau nhỏ hơn một chút và cho gia vị từ từ thì món ăn của hai chị sẽ ngon hơn nhiều ạ.''

Tôi gật đầu đồng ý, mỉm cười với Nene: ``Lần sau tôi sẽ thử lại.''

Nene cười rạng rỡ, mắt sáng long lanh như vừa nghe được ý tưởng tuyệt vời nhất: ``Vậy lần sau tôi với bà làm món khác nhé!''

Tôi nhìn chiếc chảo đen đúa đang cần ngâm rửa trong bồn rửa, thở dài một tiếng: ``Món nào ít khói hơn ấy\dots''

Nene khoanh tay, nhíu mày suy nghĩ vô cùng đăm chiêu, rồi đột nhiên thốt ra một câu xanh rờn làm tôi và Suzuno đều bất ngờ: ``Chắc là\dots\ món cháo luộc nhỉ?''

Suzuno nghe xong đành quay mặt đi chỗ khác, khẽ che miệng bật cười nghiêng ngả trước sự ngốc nghếch đáng yêu của Nene, còn tôi thì không thể nào nhịn được nữa, bật cười lớn tiếng trong căn bếp vừa trải qua đủ thứ thăng trầm.

\newpage

\stpt{-Histoire 3-}
\stcap{Trận game đêm khuya}
\begin{center}
    {\caviar\fontsize{12pt}{12pt}\selectfont Nhân vật: \emp{delu} \emp{polka} \emp{kaigou2} \emp{game}}
\end{center}

Đúng khoảng chín giờ tối, chiếc điện thoại để trên bàn bỗng rung, kéo tôi ra khỏi mớ bản phối nhạc đang dần thành hình. Tôi nhặt máy lên, màn hình hiện lên tin nhắn của Polka, dòng chữ tràn đầy sự hào hứng gần như có thể bật ra khỏi màn hình:

\txtbx{``Delu-chan, tối nay bà có muốn chơi game kinh dị với tôi không? Tôi đã tìm được một trò chơi cực kỳ hay ho, không thể không thử cùng bà đâu!''}

Tôi ngoảnh đầu nhìn chiếc máy tính để bàn vẫn đang sáng màn hình, cửa sổ trình biên tập nhạc còn mở dang dở, rồi lại đảo mắt qua cánh cửa phòng thu đang khép hờ. Suy nghĩ lướt qua lịch trình của ngày mai: tôi không có buổi thu âm nào dậy sớm, có thể ngủ muộn một chút cũng chẳng sao. Sau một hồi đắn đo ngắn ngủi, tôi gõ dòng chữ đồng ý, đặt điện thoại xuống và chờ đợi.

Chẳng mất bao lâu, tiếng gõ cửa nhẹ nhàng vang lên. Khi tôi mở cửa, Polka đứng đó với một chiếc túi nhựa to đùng đựng đủ loại snack đồ ăn vặt trong tay, chiếc laptop cá nhân bịt vỏ đầy màu sắc còn được kẹp dưới nách.

Bà ấy lập tức bước vào nhà, hào hứng khoe trước mặt tôi rằng đã chuẩn bị hết mọi thứ: một trò chơi nhập vai góc nhìn thứ nhất với cốt truyện rùng rợn, đồ họa chân thực đến từng chi tiết, bảo đảm sẽ mang lại những phút giây kịch tính khó quên.

Nhìn vẻ phấn khích không thể che giấu của bà ấy, tôi vẫn không khỏi có chút e ngại. Trong lúc bà ấy cắm dây kết nối laptop với màn hình lớn của phòng thu, tôi liền hỏi nhỏ: ``Trò chơi này có quá đáng sợ không? Tôi không giỏi chịu mấy cảnh jumpscare đâu đấy.''

Nghe vậy, Polka liền vỗ mạnh vào ngực mình, ánh mắt trong đôi mắt vàng lấp lánh sự tự tin rực lửa, giọng nói đầy quả quyết: ``Yên tâm đi! Tôi là nghệ sĩ xiếc chuyên nghiệp mà, đã trải qua vô số cảnh tượng bất ngờ trên sân khấu, trên đời này chẳng có gì làm tôi sợ được đâu! Cứ để tôi dẫn dắt bà qua mọi thử thách!''

Đúng lúc bà ấy vừa dứt lời, Game - chiếc đồng hồ thông minh đang bận rộn điều chỉnh hệ thống âm thanh vòm của phòng thu - bỗng cất lên giọng nói đều đều, vô cảm như mọi khi: ``Tôi đã ghi nhận câu tuyên bố ấy vào bộ nhớ lưu trữ dài hạn. Sẽ sử dụng để đối chiếu phản ứng trong quá trình chơi game.''

Polka lập tức quay phắt sang phía chiếc đồng hồ, đôi mắt hơi nheo lại thành một đường lườm, giọng nói có chút giận dỗi: ``Ông đừng có lưu lại mấy câu thoại không quan trọng như thế chứ! Đó chỉ là sự thật thôi, không phải để ông ghi lại làm kỷ niệm đâu!''

Vừa lúc chúng tôi chuẩn bị kéo ghế ngồi xuống, tiếng bước chân chậm rãi vang lên từ hành lang bên ngoài. Chị Kaigou bước ngang qua cửa phòng, dừng chân một lúc để nhìn vào trong, không quên để lại một lời nhắc nhở đầy uy quyền, giọng nói mang vẻ nghiêm nghị: ``Hai đứa chơi thì nhớ giữ trật tự một chút, tầng trên có Ryuu với Suzu đang ngủ sớm đấy. Đừng để tiếng động quá lớn làm phiền người khác. Với lại, tuyệt đối đừng để rơi vụn bánh snack ra sàn nhà, lần trước mấy đứa để lại cả đống mảnh vỏ bánh quy, chị phải dọn mất nửa tiếng mới sạch.''

Polka nhanh nhẹn giơ ngón cái lên ra hiệu đã rõ mọi yêu cầu, gương mặt nở nụ cười tươi rói như đứa trẻ đang xin phép được chơi game. Sau khi chị Tổng bước đi xa, chúng tôi cùng nhau kéo rèm cửa phòng thu lại, hạ bớt cường độ ánh đèn từ ấm vàng sang mờ tối, tạo nên bầu không khí bí ẩn vừa đủ để đắm chìm vào trò chơi. Rồi cả hai ngồi đối diện nhau trước màn hình lớn, sẵn sàng cho một đêm đầy những bất ngờ chờ đợi phía trước.

\ct

Trò chơi bắt đầu trong khung cảnh một căn biệt thự cổ bị bỏ hoang đã lâu, bóng tối u ám bao trùm mọi góc phòng. Nhân vật mà tôi đang điều khiển trên màn hình chỉ được trang bị hai vật dụng duy nhất: một chiếc đèn pin ánh sáng yếu ớt, lập lòe như sắp tắt, và một chiếc chìa khóa cũ đã bị gỉ sét phủ lớp vàng nâu xỉn màu.

Mặc dù Polka mới chính là người kéo tôi vào cuộc phiêu lưu này, bà ấy lại liền kéo chiếc ghế gỗ sát đến sau lưng tôi, tự phong cho mình vai trò ``đạo diễn chỉ đạo chiến thuật'' với nhiệt tình tràn đầy không thể che giấu. Tôi cũng tỏ vẻ nghi ngờ khi bà ấy lại cầm tay chỉ việc thay vì trực tiếp chơi, có lẽ nào\dots

Polka vừa nhai nhóp nhép gói snack khoai tây sần sật vừa ghé sát tai tôi, giọng thì thầm gấp gáp: ``Tôi nghĩ bà nên điều khiển nhân vật đi kiểm tra căn phòng nằm bên trái cánh cửa kia trước đã.''

Tôi ngẩn người nhìn cánh cửa gỗ mun mốc, u ám trước mặt nhân vật trong game rồi ngập ngừng đáp: ``Nhưng cánh cửa bên trái đó đang khóa chặt mà.''

``Thì bà nhấn nút tương tác xem sao, biết đâu mở được!'' bà ấy giục giã không chờ đợi.

Tôi đành làm theo lời Polka, nhấn mạnh phím E trên bàn phím để kích hoạt lệnh tương tác. Ngay khoảnh khắc cánh cửa gỗ vừa nhích ra một khe hẹp, một tiếng động lớn chát chúa, như có gì đó đập mạnh vào tường, bỗng dưng vang lên từ bóng tối sâu thẳm bên trong. Đáng nói là dù không phải người trực tiếp cầm chuột điều khiển, Polka lại là người đầu tiên cất tiếng hét thất thanh lên, tiếng hét vang lên đột ngột đến mức làm tôi cũng giật mình thót tim, tay buông thả ra khỏi bàn phím khiến nhân vật trong game lảo đảo đập sầm mặt vào bức tường gạch ngay trước mắt.

Thấy tôi đứng hình chết lặng trước màn hình, Polka vội vàng tìm cách chữa thẹn, giọng nói bập bẹ như vừa bị bắt quả tang: ``Tôi\dots\ tôi chỉ đang thử âm lượng giọng mình trong phòng xem có chuẩn không thôi!''

Tôi nuốt nước bọt đánh ực, tay chỉ thẳng vào góc màn hình nơi vừa có bóng đen vụt qua: ``Nhưng trong game vừa có cái gì đó chạy qua thật kìa!''

Polka vội gật gù cái đầu, cố gắng tỏ ra bình thản như mọi chuyện đã nằm trong dự tính của bà: ``Đúng rồi, tôi bảo bà mở cửa là để xác nhận điều đó đấy!''

Ngay lúc ấy, chiếc đồng hồ thông minh Game bỗng cất tiếng cảnh báo đều đặn, thông báo: ``Tiếng ồn của Omaru-san vượt quá ngưỡng an toàn cho phép của Kaigou-user,'' khiến Polka lật đật bịt miệng, hạ giọng thì thầm vào tai tôi: ``Được rồi, từ giờ tôi hứa sẽ giữ im lặng tuyệt đối, không phát ra tiếng động nào nữa.''

Tôi chỉ đáp lại với một cái thở dài, chẳng còn hơi sức nào để phản bác lại lời bào chữa quá đỗi vô lý của bà ấy.

\ct

Trước mắt nhân vật tôi điều khiển bây giờ là một hành lang dài hun hút, mờ ảo trong bóng tối, và ở cuối con đường đó là một cánh cửa gỗ khác phủ đầy những vết cào xé rùng rợn, như có thứ gì đó đã cố gắng xé tung nó ra.

Dù trong lòng tôi chẳng muốn tiến lại gần nơi đó chút nào, Polka vẫn không ngừng xúi giục từ phía sau lưng: ``Không sao đâu, cứ bước nhẹ nhàng tiến từng chút một thôi, tôi ngồi ngay bên cạnh đây, bảo kê cho bà mà!''

Sau khi miệt mài lục tung từng ngóc ngách trong căn phòng cũ kỹ, cuối cùng tôi cũng tìm thấy chiếc chìa khóa thứ hai ẩn sâu trong ngăn kéo của cái tủ gỗ cũ mòn. Từ đó, nhân vật mà tôi điều khiển đảm nhận vai trò dẫn đường, giữ chiếc đèn pin yếu ớt soi từng bước chân trong bóng tối, còn Polka thì dồn toàn bộ sự tập trung, ghé sát màn hình đến mức trán gần như chạm vào kính, cố gắng giải mã từng dòng chữ nguệch ngoạc ghi trên mẩu giấy cũ màu vàng mà nhân vật vừa nhặt được giữa đống đồ đạc bỏ hoang.

Bà ấy vừa đọc được nội dung, liền nhanh chóng đọc lại nhỏ giọng, hơi thở thì thầm đầy hồi hộp: ``Nó ghi rằng cánh cửa gỗ ở cuối hành lang dài kia chỉ có thể mở được khi mình bật lại cầu dao điện nằm sâu dưới tầng hầm của căn biệt thự này.''

Không có lựa chọn nào khác, chúng tôi đành phải điều khiển nhân vật quay bước, lùi lại từ đầu hành lang để đi tìm lối xuống tầng hầm tăm tối, nơi bóng tối dày đặc như có thể nuốt chửng bất cứ ai bước vào. Mỗi bước chân nhân vật đặt trên bậc thang gỗ mục nát đều phát ra tiếng cót két rợn người, tiếng vang vọng trong không gian tĩnh lặng càng làm tăng thêm sự hồi hộp khó tả.

Khi tôi đột nhiên dừng nhân vật lại ngay ở bậc thang thứ ba, vì cảm giác lạnh sống lưng rằng có điều gì đó bất thường đang ẩn mình trong bóng tối phía sau, Polka cũng lập tức nín thở theo, cơ thể cứng đờ rồi nhỏ giọng thì thầm hỏi tôi lý do dừng lại.

``Tôi nghe thấy tiếng bước chân dồn dập, nặng nề đang tiến lại ngay sau lưng nhân vật\dots'' tôi đáp nhỏ, giọng run run. Trong không gian tĩnh mịch đến đáng sợ của trò chơi, một tiếng bước chân chậm rãi, đầy đe dọa quả thực đang tiến lại gần, mỗi bước đều như đập vào tim của cả hai.

Polka cố gắng hít một hơi sâu để trấn tĩnh bản thân, rồi đập nhẹ vào vai tôi với giọng thì thầm nhưng đầy gấp gáp: ``Tôi nghĩ bà nên nhấn nút chạy nhanh đi! Không thì nó bắt kịp nhân vật mất!''

``Bà vừa mới hứa là sẽ không phát ra tiếng động lớn nữa mà?'' tôi nhắc khéo, mắt vẫn không rời khỏi màn hình.

``Thì xúi bà chạy đâu có tính là hét! Đây là chỉ dẫn chiến thuật mà!'' cô bạn Cáo sa mạc nhanh chóng bào chữa, vẫn giữ giọng nhỏ nhất có thể.

Nhưng tất cả sự bình tĩnh vứt đi ngay khoảnh khắc một bóng đen u ám, to lớn vụt qua ngay sau lưng nhân vật trên màn hình. Polka không kìm được sự hoảng loạn, liền cất tiếng hét toáng lên, tiếng hét làm tôi cũng bất giác hét theo, tay múa may ấn lung tung các phím trên bàn phím để cố gắng điều khiển nhân vật chạy trốn.

Ngay lập tức, hệ thống thông minh của Game tự động can thiệp, hạ âm lượng của loa ngoài xuống mức thấp nhất, như một biện pháp ngăn chặn kịp thời trước khi có thêm tiếng ồn nào nữa vang ra ngoài phòng thu.

Vừa kịp bình tĩnh lại sau cơn hoảng loạn khi thấy bóng đen đuổi theo nhân vật, tiếng hét của cả hai vẫn còn vương trên môi thì từ phía cầu thang ngoài hành lang bỗng vang lên tiếng bước chân giậm mạnh như sấm, từng nhịp đập đầy sát khí lan tỏa khắp cả tầng nhà.  Cả tôi và Polka lập tức cứng người, im bặt không dám thở mạnh, như thể chỉ cần một hơi thở nhỏ cũng sẽ làm bùng lên cơn thịnh nộ của người đang đứng trên kia.

Từ tầng ba, giọng nói lạnh ngắt như băng của chị Kaigou vọng xuống, xé tan bầu không khí hồi hộp của trò chơi: ``Đã dặn hai đứa chơi nhỏ tiếng rồi cơ mà!''

Cô bạn Cáo sa mạc giật mình đến mức ôm chầm lấy chiếc laptop vào ngực, như thể nó là chiếc phao cứu sinh duy nhất giữa biển lửa, rồi vội vàng gọi vọng ra ngoài: ``Xin lỗi, Kaigou-senpai! Tại trong game có cảnh hù dọa bất ngờ quá ạ!''

``Cảnh bất ngờ kiểu gì mà làm rung chuyển cả ba tầng nhà thế hả?'' giọng chị Kaigou đáp lại đầy đe dọa, làm cho cả hai tôi đều lặng người đi.

Tôi quay mặt nhìn vào màn hình game đã tự động tạm dừng, nơi nhân vật tôi điều khiển đang đứng đối mặt với một con búp bê gỗ cũ kỹ, thân thể nó xước xát đầy những vết kỳ lạ, trông như vừa bước ra từ trong truyện ma. Ngay khoảnh khắc ấy, cái đầu gỗ tròn của con búp bê đột ngột xoay ngược 180 độ, mặt nó hướng thẳng vào màn hình, đôi mắt thủy tinh trông như đang nhìn thẳng vào chúng tôi. Cả tôi và bà đều nín thở, cùng nhau giơ tay bịt chặt miệng để không phát ra thêm tiếng động nào.

Đúng lúc tim cả hai đang đập thình thịch trong ngực, một chiếc dép nhãn dán đơn màu đen từ tầng ba bất ngờ lao xuống cầu thang với tốc độ đáng kinh ngạc, lướt qua cánh cửa phòng tầng hai chúng tôi đang ngồi rồi đáp xuống tấm thảm dày ngay trước ngưỡng cửa một cách gọn gàng, như một phi cơ hạ cánh chính xác.

Chúng tôi nhìn chiếc dép nằm im lìm trên sàn, rồi quay sang nhìn nhau trong câm nín, không ai dám nói một lời. Chị Tổng lại vọng xuống lần nữa, sắc lệnh không thể tranh cãi: ``Giảm âm lượng. Ngay lập tức.''

Polka lặng lẽ cúi người, bò từng bước nhỏ ra cửa, nhặt chiếc dép lên và đặt ngay ngắn lại bên ngoài bậc thang, rồi quay lưng vào trong, thì thầm nhỏ đến mức chỉ có tôi nghe được: ``\hl{Bà có muốn tiếp tục chơi không?}''

Sau cú giật mình đôi, tôi lắc đầu lia lịa, ra hiệu với bà rằng phải chỉnh âm lượng trong game xuống mức tối thiểu, và tuyệt đối không được cất ra bất cứ tiếng động nào nữa, nếu không muốn có một chiếc dép khác bay về phía mình.

\ct

Sáng hôm sau, khi cả nhà đã tập trung đầy đủ trong phòng khách rộng rãi, Polka hào hứng đứng giữa, tay vẫy vẫy như đang diễn thuyết trước một đám đông, hớn hở kể lại sự cố chiếc dép bay qua các tầng nhà tối hôm qua cho mọi người cùng nghe, giọng nói còn đầy ắp sự ngỡ ngàng không thể che giấu.

``Mọi người có tin không? Tối qua chiếc dép của Kaigou-senpai bay từ tầng ba xuống phòng chúng tôi nhanh hơn cả tên lửa! Em còn nghĩ có tên lửa đạn đạo lao qua cửa nữa là!'' Cô ấy vỗ tay đánh đập, mặt rạng rỡ như vừa chứng kiến một kỳ tích thế giới, khiến Suzuno ngồi bên cạnh cũng phải bụm miệng cười, còn Ryuuhou ngồi cạnh lắc đầu cười không ngừng.

Chị Kaigou thì thong dong ngồi trên ghế sofa bọc vải xám, hai tay khoanh chặt trước ngực, vẻ mặt vẫn giữ nguyên vẻ nghiêm nghị cứng nhắc như mọi khi, rồi từ từ nhả ra từng câu trả lời điềm tĩnh: ``Chị không hề có ý định ném dép vào ai cả đâu. Đó chỉ là chị vứt nhẹ xuống sàn để cảnh báo hai đứa thôi, chứ có ý gì đâu.''

Nghe vậy, Polka lập tức nhảy cẫng lên, giơ ngón tay chỉ thẳng về phía cầu thang gỗ ở đầu phòng, mặt hờn dỗi vẻ không phục chút nào: ``Nhẹ đâu mà nhẹ! Chiếc dép đó bay thẳng một mạch từ tầng ba xuống tận cửa phòng tầng hai của bọn em, đó gọi là vứt nhẹ sao ạ?''

Chị Tổng vẫn không thay đổi sắc mặt, thản nhiên liếc nhìn cô bạn Cáo sa mạc rồi đáp một câu xanh rờn: ``Chị kiểm soát được đường bay của mình, đương nhiên nó sẽ rơi đúng vị trí chị muốn rồi.''

Đúng lúc đó, chiếc đồng hồ thông minh Game đang đặt trên kệ tủ đột nhiên chen ngang vào câu chuyện, giọng nói đều đều, vô cảm như mọi khi nhưng lại mang đậm giọng điệu phân tích kỹ thuật chuyên nghiệp: ``Sau khi tính toán lại góc phóng, quãng đường bay và điểm rơi cuối cùng của vật thể, chiếc dép đã được phóng đi với độ chính xác lên đến 98\%, cùng khả năng tận dụng lực cản không khí để duy trì quỹ đạo ổn định, những con số này đạt tiêu chuẩn đáng kinh ngạc so với một hành động phi chuyên nghiệp.''

Chị Kaigou lập tức quay phắt sang nhìn chiếc đồng hồ thông minh, ánh mắt lườm gai gai như muốn đốt cháy nó, giọng nói đầy chất vấn: ``Này Game, ông đang đứng về phía ai đấy hả? Lại còn đem ra phân tích nữa!''

``Tôi chỉ đứng về phía dữ liệu thực tế được thu thập được thôi,'' ông ta đáp trả không chớp mắt, giọng nói vẫn giữ nguyên sự thản nhiên không hề nao núng trước ánh mắt giận dữ của chị Tổng.

Polka tranh thủ đúng lúc chị ấy đang dồn toàn bộ sự chú ý vào việc lườm chiếc đồng hồ thông minh, liền rón rén dịch chuyển ghế đến sát bên tôi, ghé tai tôi thì thầm nhỏ: ``\hl{Lần sau chúng mình chuyển sang chơi game giải đố nhẹ nhàng thôi nhé, đừng dại dột mà đụng đến game kinh dị nữa, không lại có thêm một chiếc dép khác bay về phía đầu thì chết.}''

Mặc dù tôi chỉ biết gật đầu lia lịa đồng ý với đề xuất hợp lý của Polka, thì chị Kaigou ở phía sau vẫn tinh ý nghe thấy toàn bộ câu nói, liền lạnh lùng chêm vào một câu như một sắc lệnh không thể tranh cãi: ``Lần sau mà muốn chơi gì thì cũng phải chọn ban ngày mà chơi, đừng có mà đêm hôm khuya khoắt lại gây ồn ào cả nhà.''

\ct

Đêm ấy, chúng tôi vẫn tiếp tục vật lộn với màn hình game trong không khí yên lặng đến ngột ngạt. Dù không còn những tiếng hét thất thanh vang lên làm rung chuyển cả căn phòng như ban đầu, Polka vẫn không ngừng giật mình thon thót mỗi khi trò chơi bật ra một tiếng động lạ từ bóng tối, còn tôi thì siết chặt con chuột máy tính đến mức các ngón tay cứng đờ, cổ tay mỏi rần chỉ sau vài tiếng đồng hồ cầm nắm.

Nhưng bù lại cho những lần giật mình liên tục, chúng tôi lại bật cười nghiêng ngả mỗi khi tôi vô tình điều khiển nhân vật chạy lạc hướng, hay đâm sầm đầu vào bức tường gạch trong game khiến màn hình rung lên bần bật. Đôi khi chỉ vì một lỗi điều khiển nhỏ, cả hai lại phải che miệng cười đến mức bụng đau, sợ rằng tiếng cười sẽ vô tình vượt ra ngoài phòng thu và thu hút thêm cơn thịnh nộ của chị Kaigou. Đến gần nửa đêm, sau bao nhiêu khó khăn và hồi hộp, nhân vật tôi điều khiển cuối cùng cũng tìm được cánh cửa thoát hiểm cuối cùng ra khỏi căn biệt thự ma quái ám ảnh suốt đêm.

Khi dòng chữ lớn màu vàng ``CHIẾN THẮNG'' hiện lên đầy màn hình, Polka mừng rỡ giơ hai tay lên trời như vừa đoạt được giải thưởng lớn, nhưng kịp nhớ ra lời cảnh báo nghiêm khắc của chị Kaigou, bà chỉ dám reo lên một tiếng nhỏ xíu, vừa đủ để không quấy rầy tầng trên: ``Chúng ta thắng rồi!''

Tôi mỉm cười nhìn bà, giọng cũng chỉ đủ để hai người nghe: ``Hồi nãy bà lại trót reo lên rồi đấy nhé.''

``Tôi chỉ phát ra đúng âm lượng của sự vui mừng thôi mà!'' Polka vội vàng bào chữa, giọng thì thầm như sợ ai nghe thấy.

Tôi quay đầu nhìn ra ngoài cánh cửa phòng thu, chiếc dép đen của Kaigou vẫn nằm im lìm, ngay ngắn trên tấm thảm dày ở bậc thang. Polka sau đó cẩn thận gấp gọn những gói snack còn dở dang, bỏ từng thứ một vào chiếc balo màu sắc của mình. Trước khi bước ra cửa, bà còn quay lại hứa hẹn với tôi, lần sau nhất định sẽ chọn một trò chơi nhẹ nhàng, ít cảnh hù dọa hơn nhiều, không để xảy ra cảnh tượng bị ném dép như đêm nay nữa.

Khi tôi mỉm cười hỏi liệu bà có thực sự chắc chắn về lời hứa đó không, Polka lập tức ngẩng đầu lên với nụ cười rạng rỡ như hoa nở, đáp lại ngay lập tức: ``Chắc chắn luôn! Nhưng tôi không hứa là mình sẽ không reo lên đâu đấy nhé!''

\newpage

\stpt{-Histoire 4-}
\stcap{Lời khuyên của sư tử}
\begin{center}
    {\caviar\fontsize{12pt}{12pt}\selectfont Nhân vật: \emp{delu} \emp{botan}}
\end{center}

Vào một buổi sáng cuối tuần yên ả, khi tôi vẫn đang mải miết sắp xếp các bản ghi âm cũ trong phòng thu, điện thoại bỗng rung lên. Tin nhắn từ Botan hiện lên, lời văn ấm áp như một lời mời gọi nhẹ nhàng: ``Delu-chan, bà có muốn cùng tôi thử một buổi tập nhẹ ở phòng gym không? Chỉ vận động cho cơ thể thoải mái thôi, không cần phải cố gắng quá sức đâu.''

Tôi ngước nhìn ra cửa sổ, ánh nắng vàng nhạt len lỏi qua kẽ lá, rồi lại cúi xuống nhìn đôi giày thể thao còn nguyên tem mác đặt dưới gầm bàn. Đôi giày đó tôi đã hào hứng mua từ tháng trước, với bao dự định sẽ tập luyện đều đặn, thế mà mãi đến nay vẫn chưa một lần được xỏ chân vào. Những công việc ở phòng thu luôn chất đống, khiến tôi cứ dời ngày này sang ngày khác. Bây giờ, lời đề nghị của bà như một cái cớ hoàn hảo để tôi tạm gác lại những lo toan thường nhật. Sau vài giây đắn đo, tôi gõ dòng chữ đồng ý, rồi bắt đầu chuẩn bị quần áo tập luyện.

Chẳng mấy chốc, cô bạn Sư tử đã hiện ra ở sảnh chung của căn nhà, nụ cười quen thuộc nở trên môi khi thấy tôi bước xuống. Bà mặc bộ đồ tập thể thao đơn giản, dáng người cao lớn, vững chãi như có thể che chắn cho mọi thứ xung quanh, chiếc khăn bông nhỏ vắt trên vai càng tôn lên vẻ năng động.

``Đúng là tôi còn phải làm quen lại với việc vận động hằng ngày lắm,'' tôi vừa nói vừa cúi xuống gắn lại dây giày thể thao, giọng hơi ngượng ngùng.

Botan lập tức tiến lại gần, đặt tay nhẹ lên vai tôi, giọng ấm áp trấn an: ``Không sao cả, lần đầu tiên ai cũng thế thôi. Tôi cũng từng mất một thời gian dài để làm quen với lịch trình tập luyện đều đặn mà. Tới đây chúng ta cùng nhau điều chỉnh, không vội vàng gì cả.''

\ct

Chúng tôi đến một phòng gym yên tĩnh nằm không xa nhà lắm, nơi đó không đông đúc ồn ào như những trung tâm thể thao lớn, mà mang lại cảm giác thư thái, thoáng đãng với những ô cửa sổ lớn tràn ngập ánh sáng.

Botan rất tỉ mỉ, dắt tôi đi qua từng động tác khởi động nhẹ nhàng, không quên ôn tồn nhắc nhở từng chút một: ``Chúng mình đến đây để chăm sóc sức khỏe chứ phải đi thi đấu đâu. Bất cứ lúc nào bà cảm thấy mệt mỏi hay hơi ngợp, cứ dừng lại ngay lập tức, không cần cố gắng chịu đựng.''

Tôi ngoan ngoãn gật đầu làm theo. Nhưng cơ thể dường như đã quá lâu không vận động, những động tác giãn cơ đơn giản cũng khiến bờ vai tôi cứng nhắc, lúng túng không biết đặt tay chân vào đâu. Thấy vậy, Botan liền tiến lại gần, đôi bàn tay ấm áp nhẹ nhàng nắm lấy vai tôi, chỉnh sửa tư thế từ từ: ``Đừng gồng cổ quá đấy bà. Hít thở sâu vào, thả lỏng toàn bộ cơ thể ra, rồi thực hiện mọi thứ thật chậm thôi.''

Chúng tôi cứ thế thực hiện những bài tập đơn giản, xen kẽ những khoảng nghỉ ngơi ngắn giữa từng lượt, không hề đặt ra bất kỳ áp lực nào về trọng lượng tạ hay quãng đường chạy trên máy.

Thời gian trôi qua lặng lẽ, cho đến khi hơi thở của tôi bắt đầu dồn dập, mồ hôi rịn ra trên trán. Botan lập tức giơ tay ra hiệu cho tôi dừng lại, rồi chủ động lấy chai nước từ tủ lạnh trao cho tôi, trước khi ngồi xuống chiếc ghế dài dài bên cạnh tôi để cùng nghỉ ngơi.

\ct

Tôi vừa dùng ống tay áo quệt đi những hạt mồ hôi đang đọng trên trán, mắt vừa thở dốc vừa ngước nhìn những người tập gym xung quanh, có người đang dồn lực nâng tạ nặng đến mặt đỏ bừng, có người lại chạy bộ trên máy với tốc độ không ngừng nghỉ. Lấy lại một chút hơi thở ổn định, tôi quay sang nhìn người bạn sư tử bên cạnh, khẽ cười nói: ``Tôi cứ nghĩ một người khỏe mạnh, dũng mãnh như bà sẽ chọn những bài tập nặng đô, hấp lực hơn cơ, chứ không ngờ hôm nay lại đi cùng tôi tập nhẹ nhàng thế này.''

Botan mỉm cười, đôi mắt nâu ấm áp sáng lên như có ánh nắng mặt trời chiếu vào, ánh nhìn kiên định mà dịu dàng, không chút nào vội vàng: ``Hôm nay tôi đến đây là để đồng hành, tập cùng bà chứ đâu phải để một mình tôi cố gắng đâu. Tôi không muốn bắt bà phải gồng mình chạy theo tốc độ của mình, chỉ muốn chúng ta cùng nhau vận động thoải mái thôi.''

Tôi cúi đầu xuống, nhìn xăm soi hai bàn tay của mình vẫn còn hơi run rẩy sau bao nhiêu động tác tập luyện vừa rồi, rồi thở một hơi dài nặng nề: ``Tôi chỉ muốn bản thân mình có đủ sức khỏe để làm được mọi thứ thôi\dots\ Từ sáng tác nhạc, dồn thời gian vào các buổi thu âm trong phòng cho đến việc gặp gỡ, đi chơi cùng tất cả mọi người. Mỗi lần tôi phải dừng lại nghỉ giữa chừng, không thể tiếp tục được nữa, tôi lại có cảm giác như mình đang tự lùi lại phía sau, bỏ rơi tất cả những người đang đi cùng mình.''

Botan không vội trả lời câu nói của tôi. Bà ấy chậm rãi vặn mở nắp chai nước khoáng đang đặt trên ghế, uống một ngụm nhỏ để làm ấm cổ họng rồi mới quay sang nhìn tôi. Trong ánh mắt của người bạn sư tử vốn luôn mang vẻ kiên cường, vững chãi như núi non ấy, bỗng dưng thoáng qua một gợn sóng trầm mặc, như có một câu chuyện lâu ngày đang chìm sâu trong đáy lòng.

Cô bạn Sư tử đặt bàn tay to lớn, ấm áp của mình nhẹ nhàng lên vai tôi, như một điểm tựa vững chãi đến kỳ lạ mà chỉ cần một cái chạm ấy, mọi mệt mỏi trong người tôi đều tan biến đi phần nào.

``Không cần phải bắt bản thân mình lúc nào cũng phải mạnh mẽ đâu, Delu-chan. Điều quan trọng nhất là bà biết khi nào nên đứng dậy tiếp tục bước đi, và khi nào thì cần cho phép bản thân mình nghỉ ngơi để lấy lại sức mà đi tiếp,'' cô bạn Sư tử thả nhẹ giọng, như đang chia sẻ một sự thật mềm mại mà bà đã mất nhiều thời gian để hiểu thấu, ``cơ thể của bà không phải là đối thủ để bà chiến đấu hay phải đánh bại nó đâu. Nó là người bạn đồng hành duy nhất, sẽ đi cùng bà qua tất cả những ngày tháng dài phía trước. Nếu bà cứ coi mỗi lần mệt mỏi, mỗi lần phải dừng lại là một thất bại, thì bà sẽ chẳng bao giờ cho bản thân mình có cơ hội để phục hồi, để mạnh mẽ trở lại đâu.''

Tôi ngước lên nhìn Botan, cảm nhận được sự bao bọc, che chở sâu sắc trong từng câu chữ bà nói ra, rồi khẽ nói: ``Bà nói như thể\dots\ bà đã từng phải trải qua và học được bài học đó một cách rất đau đớn vậy.''

Botan lặng đi một nhịp, nụ cười nhẹ nhàng thường ngày trên môi bà bỗng thoáng qua một nét trầm lắng khó tả. Bà nhìn thẳng về phía trước, vào khoảng không gian rộng lớn của phòng gym, ánh mắt lộ rõ sự quyết tâm và một lời hứa âm thầm đang gieo mầm từ trong đáy mắt.

``Tôi nghĩ ai rồi cũng phải học bài học đó thôi\dots\ Ngày xưa, tôi từng có một người bạn rất quan trọng với mình. Tôi đã luôn muốn đứng ra che chắn cho cô ấy, muốn dốc hết tất cả sức lực của mình để bảo vệ cô ấy đến cùng trước mọi sóng gió, mọi khó khăn mà cô ấy phải đối mặt. Nhưng rốt cuộc\dots\ lời hứa ấy lại đành phải dở dang. Chuyện đó vẫn luôn là một vết thương chưa lành, cứ nằm yên trong lòng tôi đến tận bây giờ.''

Botan quay sang nhìn thẳng vào mắt tôi, giọng nói của bà trở nên đanh thép, dũng mãnh như một lời thề, nhưng lại đong đầy sự chân thành tuyệt đối, không chút nào giả vờ: ``Vì vậy, đứng trước bà lúc này, tôi không muốn lịch sử đó lặp lại lần nữa. Tôi sẽ không để bà phải một mình gồng gánh mọi thứ, hay phải kiệt sức mà rơi lại phía sau một mình. Cứ yên tâm mà bước tiếp trên con đường bà chọn, bất cứ khi nào bà cảm thấy mệt mỏi, hay muốn gục ngã vì quá nhiều áp lực, tôi sẽ luôn ở đây, làm lá chắn cho bà, và bảo vệ bà đến cùng.''

Lời khẳng định đong đầy sức mạnh và sự chân thành của Botan khiến lồng ngực tôi dâng lên một cảm giác ấm áp đến lạ kỳ. Nhìn thấy ánh mắt bao bọc dũng mãnh ấy, những lo âu và áp lực đè nặng trong tôi như dần tan biến. Chúng tôi đứng dậy, quay lại thực hiện thêm một lượt tập ngắn trong sự thư thái hoàn toàn.

\ct

Khi buổi tập kết thúc, tuy tôi không làm được điều gì quá phi thường hay vượt ngưỡng, nhưng tâm trí và cơ thể tôi lại cảm thấy nhẹ nhõm, dễ chịu hơn rất nhiều so với lúc mới bước vào. Botan mỉm cười khoác vai tôi đi ra cửa phòng gym, giọng ấm áp như những tia nắng cuối ngày: ``Lần này cảm thấy thế nào, Delu-san? Có thấy cơ thể mình nhẹ nhàng hơn không?''

Tôi khẽ rung rung hai vai, cảm nhận sự thư thái lan tỏa từ từng thớ cơ, rồi mỉm cười đáp: ``Ừ\dots\ thật kỳ lạ sao ấy. Lúc đầu tôi còn sợ sẽ mệt lả không đi nổi về, mà giờ lại thấy có thêm sức mạnh dồi dào thế nào ấy.''

Botan nở một nụ cười ấm áp, đôi vai lớn vững chãi của cô bạn Sư tử khẽ rung lên theo tiếng cười. Chị nghiêng người về phía tôi một chút, bàn tay vỗ nhẹ nhàng hai lần trên vai tôi như một cái vỗ về an ủi, sau đó giữ nguyên đó một lúc để truyền cho tôi sự ấm áp và chắc chắn.

Giọng cô ấy trầm ấm, chứa đựng sự chân thành vô bờ bến khi nhắn nhủ từng lời, như thể đang dặn dò một người bạn cùng chí cốt: ``Lần sau nếu vẫn muốn cùng tôi đến phòng gym tập luyện, hay thậm chí chỉ đơn giản là có chuyện gì đó muốn tâm sự, muốn tìm một chỗ để dựa dẫm thì cứ gửi tin nhắn cho tôi bất cứ lúc nào nhé. Không bao giờ được ngại ngần hay nghĩ mình là người phiền phức đâu, được chứ? Bọn mình là bạn bè thân thiết mà, tất cả mọi chuyện, dù vui hay buồn, chúng ta đều phải cùng nhau nắm tay bước qua.''

Tôi mỉm cười gật đầu, giọng nhỏ đầy biết ơn: ``Cảm ơn bà nhiều lắm, Botan-chan. Nhờ bà mà có lẽ\dots\ tôi cũng thanh thản tâm hồn hơn.''

Lần này, tôi không vội vã hứa hẹn sẽ tập luyện thật chăm chỉ hay trở nên phi thường nữa, mà tôi thầm hứa sẽ học cách lắng nghe cơ thể mình, và trân trọng sự bảo vệ đầy ấm áp mà Botan đã dành cho tôi.

\newpage

\stpt{-Histoire 5-}
\stcap{Buổi karaoke holoFive}
\begin{center}
    {\caviar\fontsize{12pt}{12pt}\selectfont Nhân vật: \emp{delu} \emp{nene} \emp{lamy} \emp{polka} \emp{botan}}
\end{center}

Vài ngày sau, điện thoại tôi bỗng rung liên tục, thông báo tin nhắn mới từ nhóm chat chung của cả năm đứa. Tôi mở màn hình lên, thấy tin nhắn đầu tiên từ Nene nhấn nhá đầy hào hứng: ``Cuối tuần này chúng mình rảnh cả, cùng tổ chức một buổi karaoke nhé bà con ơi!''

Chẳng mấy chốc, Polka đã nhảy vào trả lời bằng một chuỗi dài biểu tượng pháo hoa lấp lánh, đủ thấy bà ấy phấn khích đến mức nào. Lamy cũng nhanh chóng gửi một dấu tick xanh lớn, báo hiệu đồng ý tuyệt đối. Botan thì gửi câu hỏi ngắn gọn, thực tế như mọi khi: ``Vậy tổ chức ở đâu đây?''

Tôi nhìn căn phòng thu của mình, nơi luôn có đủ micro chuyên nghiệp, loa kiểm âm chất lượng cao, và không gian rộng rãi thoải mái để tất cả mọi người có thể ngồi lại cùng nhau. Tôi liền gõ nhanh dòng đề nghị vào nhóm: ``Hay là chúng mình tổ chức tại phòng thu của tôi đi? Ở đây đủ chỗ, âm thanh cũng tốt để hát hò.''

Nene lập tức phản hồi bằng một chuỗi trái tim đỏ rực, giọng mừng rỡ hiện rõ qua từng chữ: ``Tuyệt vời quá! Vậy chốt nhé! Delu lo chuẩn bị phòng, tôi lo toàn bộ đồ ăn thức uống cho cả nhóm!''

Tôi đọc lại tin nhắn lần nữa, mỉm cười rồi đứng dậy bắt đầu dọn dẹp căn phòng. Tôi không cần phải làm cho nó trông hoàn hảo hay cầu kỳ gì đâu, chỉ đơn giản là sắp xếp lại mọi thứ gọn gàng, để mọi người có đủ chỗ đặt cốc nước, túi xách và những món đồ mang theo.

Suzuno cũng sang giúp tôi mang thêm vài chiếc ghế mềm từ phòng khách lên, còn anh Kuon và Ryuuhou thì tỉ mỉ kiểm tra từng sợi dây nối micro, sắp xếp chúng gọn gàng dọc theo tường để không ai vô tình vấp phải khi đi lại. Game thậm chí còn kích hoạt thêm cả chế độ cách âm tuyệt đối nhằm không làm phiền tất cả mọi người xung quanh; mà quả thực, tại sao ông ta lại không làm điều đó khi Polka hú hét chứ?

Tôi còn đặt thêm một chiếc ghế nhỏ cạnh bộ mixer âm thanh, phòng khi có ai muốn nghỉ ngơi trong lúc hát, hay đơn giản là muốn ngồi ngoài một chút để thưởng thức bạn bè mình thể hiện.

\ct

Chiều thứ Bảy cuối tuần, lần lượt cả bốn người bạn của tôi đều đến đủ đầy. Nene là người đầu tiên tới, tay xách hai túi lớn nhỏ chứa đầy đồ ăn vặt, từ bim bim, kẹo ngọt đến những hộp trái cây đã được rửa sạch. Lamy thì đem theo một hộp bánh nhỏ xinh tự làm, thơm phức mùi bơ và vani, đủ cho cả năm người cùng thưởng thức.

Polka đến với chiếc túi đựng đạo cụ quen thuộc của bà ấy, nhưng vừa bước vào cửa đã hứa hẹn ngay: ``Hôm nay tôi không đem đạo cụ ra biến trò chơi karaoke thành màn ảo thuật đâu, yên tâm!'' Còn Botan thì mang theo một cuốn sổ nhỏ ghi danh sách những bài hát đã chuẩn bị kỹ lưỡng từ trước, để chúng ta không phải mất thời gian suy nghĩ.

``Ai muốn hát mở màn trước đây?'' Nene cất tiếng hỏi, tay cầm micro lắc lư đầy hào hứng. Nhưng không ai trả lời ngay, mọi người đều ngần ngại nhường nhịn nhau. Polka liền chỉ thẳng vào tôi, cười tinh nghịch: ``Chủ nhà phải mở màn chứ sao! Delu-chan hát trước đi!''

Tôi với tay cầm lấy micro, nhưng rồi lại do dự đặt nó xuống bàn, ngập ngừng nói: ``\dots Tôi có thể hát ở cuối cùng được không? Tôi muốn ngồi nghe mọi người hát trước đã.''

Nene không hề thúc ép tôi, bà ấy liền gật đầu ngay: ``Được chứ, không sao cả. Ai muốn hát trước thì cứ chọn bài thôi.''

Botan mở cuốn danh sách bài hát ra, lần lượt mọi người cùng chọn. Lamy chọn một bản nhạc chậm, du dương để làm nóng không khí. Nene chọn một bài vui nhộn, giai điệu bắt tai để mọi người cùng vỗ tay theo. Còn Polka thì yêu cầu một ca khúc có nhịp nhanh, sôi động, để tất cả cùng đứng lên nhún nhảy.

Trong suốt buổi hát, không ai chấm điểm cho ai, cũng không ai bao giờ ghi lại hay chỉ trích những lỗi nhỏ khi hát sai lời hay lạc nhịp. Mỗi người chỉ đơn giản là hát theo cách mình thích, theo cảm xúc của mình.

``Không ai giữ được nhịp thì bọn mình cùng nhảy cùng cười hết mình thôi, Nene vừa nhún vai cười lớn, ``chúng mình đến đây để giải trí mà, có phải để thi đấu đâu!

Khi Nene đột ngột quên mất một câu hát giữa chừng, mặt cô ấy ngẩn ngơ đứng yên mồ côi nhìn vào màn hình, Polka liền nhanh chóng hát đệm vào giúp bà ấy vượt qua. ``Cứ thế này là được rồi Nene-nee, chúng ta cùng đi tiếp thôi! cô ấy reo nhỏ, kéo Nene quay lại với giai điệu.

Rồi đến lúc Polka tinh nghịch đổi toàn bộ lời bài hát thành những câu đùa cợt hài hước, Botan vẫn giữ vững nhịp điệu như một cái trụ âm thanh, thậm chí còn đệm thêm tiếng vỗ tay để hỗ trợ bạn mình. ``Polka-chan, đừng đổi lời nữa kìa! Botan cười lớn nhưng không hề làm ngắt quãng bài hát.

Lamy cười đến mức phải đặt micro xuống, vai rung lên liên tục, không ngừng bật ra những câu cười lớn. ``Cả hai người này thật sự làm tôi không ngừng cười được! bà ấy lau đi nước mắt cười, còn tôi thì chỉ ngồi yên ở góc, hai tay ôm chặt cốc trà ấm, lắng nghe từng giọng hát, từng tiếng cười lan tỏa trong căn phòng nhỏ, cảm thấy lòng ấm áp lạ thường.

\ct

Sau khi tất cả mọi người đã lần lượt hoàn thành bài hát của mình, cuối cùng cũng đến lượt tôi, Nene không cần gọi tên tôi lần nào nữa. Bà ấy chỉ đơn giản là bước tới, đặt chiếc micro nhẹ nhàng ở cạnh ghế tôi đang ngồi.

Tôi nhìn chằm chằm vào chiếc micro trước mặt, lòng bỗng dưng hồi hộp khó tả. Tôi đã chuẩn bị một bài hát mới, một tác phẩm vừa viết xong vài ngày trước. Nó không có phần đệm cầu kỳ, phức tạp gì đâu, chỉ đơn giản là giai điệu tôi tự soạn trên piano, và vài câu hát còn ngập ngừng, tôi cũng chưa chắc chắn liệu mình có muốn chia sẻ chúng với mọi người hay không.

Botan đã nhận ra sự do dự trong ánh mắt của tôi, bà ấy nhẹ nhàng nói: ``Không cần phải hát nếu bà chưa sẵn sàng đâu, không ai ép buộc gì cả.''

Polka cũng nhanh chóng gật đầu phụ họa: ``Đúng đó, bọn tôi vẫn sẽ ngồi đây nghe dù bà có hát hay không!''

Lamy thì rót thêm nước ấm vào cốc trà của tôi, như một cách an ủi âm thầm. Nene không nói gì cả, bà ấy chỉ lặng lẽ kéo chiếc ghế trống bên cạnh lại gần tôi thêm một chút, để tôi cảm thấy không đơn độc.

Tôi sâu một hơi thật sâu, rồi với tay cầm lấy chiếc micro lên, giọng hơi run khi nói: ``Tôi có một bài hát mới viết, vẫn còn chưa hoàn chỉnh lắm, mọi người đừng cười nhé.''

``Vậy càng tốt,'' Nene lập tức nói, giọng ấm áp như bao giờ. ``Bọn tôi được là những người đầu tiên nghe nó lớn lên cùng bà, còn gì tuyệt hơn.''

Tôi bật phần đệm đơn giản đã chuẩn bị sẵn. Những nốt nhạc đầu tiên nhẹ nhàng vang lên trong căn phòng nhỏ, lan tỏa khắp không gian. Tôi hát chậm rãi, giữ nhịp điệu đúng theo hơi thở của mình, để từng câu chữ trôi đi tự nhiên nhất. Bài hát không kể về một ngày cụ thể, hay một câu chuyện có đầu có cuối nào cả. Nó chỉ đơn giản nói về năm màu sắc khác nhau, cùng xuất hiện trên một tấm vải trắng. Mỗi màu có một đường nét riêng, một vẻ đẹp riêng, không giống bất kỳ màu nào khác. Nhưng khi tất cả chúng đặt cạnh nhau, chúng lại hòa quyện tạo thành một hình ảnh trọn vẹn, đẹp đến lạ kỳ.

\dots

Khi những nốt nhạc cuối cùng lắng xuống, căn phòng bỗng chốc yên lặng vài giây, không có tiếng động nào vang lên. Tôi từ từ đặt chiếc micro xuống bàn, lòng bàn tay tôi bỗng thấy ẩm ướt hồi hộp.

Nene là người đầu tiên đứng dậy. Bà ấy bước nhanh tới, ôm tôi thật chặt, giọng vui mừng nhưng cũng có chút xúc động: ``Delu-chan, bài hát của bà đẹp lắm, thực sự rất đẹp.''

Lamy cũng bước tới, vòng tay lớn ôm lấy cả hai chúng tôi, hơi ấm từ bà ấy lan tỏa sang tôi. Polka giơ hai ngón tay thành hình chữ V đầy phấn khích, nhưng tôi thấy mắt bà ấy đã long lanh ngân ngấn nước. Botan chỉ vỗ nhẹ nhàng lên vai tôi, giọng bình thản nhưng đầy trân trọng: ``Bà không cần phải sửa nó hoàn chỉnh ngay trong một đêm đâu. Hãy để cho bài hát có thời gian lớn lên cùng bà, để nó được hoàn thiện một cách tự nhiên nhất.''

Tôi mỉm cười, cố gắng giữ cho giọng mình bình thường, không bị xúc động làm thay đổi: ``Tôi sẽ hoàn thiện nó sau, sớm thôi.''

``Bất cứ khi nào bà muốn, chúng tôi đều ở đây để nghe,'' Lamy nói nhẹ nhàng.

Nene cũng gật đầu liên tục, thêm vào: ``Và lần sau, chúng ta nhất định phải lại hát cùng nhau như thế này nữa, được chứ?''

Polka lập tức giơ tay lên, như một học sinh muốn phát biểu trong lớp: ``Lần sau tôi xin được chọn bài mở màn! Tôi sẽ chuẩn bị thật kỹ!''

Botan liền nhìn vào cuốn danh sách bài hát của mình, giọng đùa cợt: ``Chỉ cần bài mở màn của bà không dài hơn mười phút là được, đừng để chúng tôi phải ngồi chờ quá lâu nhé.''

``Tôi sẽ cố gắng hết sức!'' Polka cười lớn đáp lại.

Lamy cũng bật cười theo, tiếng cười của bà lan tỏa khắp căn phòng. Tôi nhìn quanh căn phòng nhỏ của mình, thấy năm chiếc cốc được đặt cạnh nhau trên chiếc bàn lớn, năm giọng hát vừa lần lượt nối tiếp nhau trong suốt buổi chiều dài này, năm người bạn mà đã từng có lúc tôi nghĩ rằng khoảng cách giữa chúng ta sẽ kéo dài mãi mãi, không thể nào xích lại gần nhau.

Nhưng giờ đây, chúng tôi không cần phải lặp lại bất kỳ lời thề nào, không cần phải hứa hão điều gì lớn lao. Chúng tôi chỉ cần đơn giản là tiếp tục hẹn nhau một lần sau nữa, tiếp tục có những buổi chiều như thế này, để những câu chuyện, những bài hát, và tình bạn của chúng ta cùng lớn lên theo thời gian.

\end{document}